\documentclass[11pt]{article}
\usepackage[T1]{fontenc}
\usepackage{lmodern}
\usepackage{amsmath,amssymb}
\usepackage{booktabs}
\usepackage[backend=biber,style=authoryear]{biblatex}
\addbibresource{refs.bib}
\usepackage[hidelinks]{hyperref}
\newcommand{\R}{\mathbb{R}}

\title{A Tour of Lemmata}
\author{You}
\date{\today}

\begin{document}
\maketitle

\section{Write, and watch the PDF}
Type anywhere in this file. A moment after you stop, Lemmata saves and
rebuilds, and the PDF on the right updates without losing your place.

Double-click any word in the PDF to jump to it here. Press Ctrl+J to go
the other way.

\section{Maths}
Put the cursor inside the equation below to see it rendered above the line.
Your own commands work too: $\R$ is defined in the preamble.
\begin{equation}
  \int_{-\infty}^{\infty} e^{-x^2}\,\mathrm{d}x = \sqrt{\pi}
\end{equation}

\section{Citations}
Press Ctrl+Shift+C to search the bibliography, or paste a DOI there to add
a new reference. This one is already in \texttt{refs.bib}:
\textcite{knuth1984}.

\section{Mistakes}
Misspelled words get a wavy line: try right-clicking this one: recieve.
If you type an unknown command like \texttt{\textbackslash foo}, the
Problems panel explains what went wrong in plain words.

\section{Everything else}
Press Ctrl+Shift+P for every command: history, search, exports and more.
This example lives in Lemmata's own folder, so feel free to break it.

\printbibliography
\end{document}
